% ══════════════════════════════════════════════════════════════
%  Chapter 20: Embedded Database & Persistence (SQLCipher)
% ══════════════════════════════════════════════════════════════
\chapter{Embedded Database \& Persistence}
\label{ch:embedded-db}

\begin{learnbox}
\begin{itemize}
  \item Integrate SQLCipher for encrypted wallet storage across both Iced and Tauri frontends
  \item Design database schemas and implement migrations using the table-recreation technique
  \item Handle thread safety with \code{OnceLock<Mutex<Connection>>} and the singleton table pattern
  \item Implement deterministic password generation so both wallet UIs share one encrypted database file
\end{itemize}
\end{learnbox}

\lettrine{B}{oth} \index{wallet}wallet applications in this project---the \index{Iced}Iced \index{wallet}wallet (Chapter~\ref{ch:wallet-iced}) and the \index{Tauri}Tauri \index{wallet}wallet (Chapter~\ref{ch:wallet-tauri})---persist user configuration and \index{wallet}wallet addresses in an encrypted SQLite database built via \index{SQLCipher}SQLCipher. This chapter covers the shared persistence layer used by both, explaining the patterns once so that neither \index{wallet}wallet chapter needs to repeat the material.

We focus on two reader-facing goals:

\begin{itemize}
  \item \textbf{Security}: sensitive values (like API keys) are stored encrypted at rest, without requiring the user to type a passphrase.
  \item \textbf{Continuity}: wallet addresses and settings survive process restarts, so the wallet list is stable and the node connection is remembered.
\end{itemize}

\begin{warnbox}[SQLCipher encryption and key strength]
SQLCipher encryption is only as strong as the password. The deterministic password generation scheme used here trades perfect forward secrecy for the convenience of sharing one database between both wallet UIs. In a production system, evaluate whether this trade-off is acceptable for your threat model.
\end{warnbox}

\begin{notebox}
Both the Iced and Tauri wallets share the same encrypted database file. Changes made in one wallet are immediately visible in the other --- there is no synchronization delay because they operate on the same SQLite file.
\end{notebox}

---

\section{What is SQLCipher?}
\label{sec:sqlcipher-intro}

\index{SQLCipher}SQLCipher is a fork of SQLite that adds transparent 256-bit AES encryption to the entire database file. From Rust, we interact with it through the \code{rusqlite} crate (with the \code{bundled-\index{SQLCipher}sqlcipher} feature flag), which compiles \index{SQLCipher}SQLCipher from source and links it statically --- no system library installation required. The API is identical to regular SQLite; the only difference is that you supply an encryption key via \code{PRAGMA key} before any other operation. This gives us encrypted-at-rest storage with zero additional infrastructure.

The database code in both apps is nearly identical. The Iced wallet stores it in a single file (\code{src/database.rs}), while the Tauri wallet organizes it into a module directory (\code{src-tauri/src/database/mod.rs} \code{+} \code{tests.rs}). The schema, key derivation, initialization sequence, and CRUD operations are the same. We will note the few differences as they arise.

To keep this chapter readable without a repository open, the complete, verbatim source files are provided in the companion code listings appendix.

% ──────────────────────────────────────────────────────────────
\section{What we store (and why)}
\label{sec:database-schema}

We keep this persistence layer intentionally small and pragmatic. We store four tables:

\begin{itemize}
  \item \textbf{\code{settings}}: \code{base\_url}, \code{api\_key} --- required so the UI can reconnect to the node without retyping configuration each time.
  \item \textbf{\code{wallet\_addresses}}: \code{address}, optional \code{label}, \code{created\_at} --- required so the wallet list is stable and selectable after restart.
  \item \textbf{\code{schema\_version}}: \code{version} --- required so the database can evolve safely over time via migrations.
  \item \textbf{\code{users}}: \code{first\_name}, \code{last\_name}, \code{profile\_picture} (BLOB) --- present to support future UI work; in the current code the persistence primitives exist (schema + migrations), even if UI integration is minimal.
\end{itemize}

\textbf{Database Schema}

\begin{minted}[fontsize=\footnotesize,linenos=false]{text}
 +------------------------------------------+
 |          SQLCipher Encrypted DB          |
 |   (shared by Iced and Tauri wallets)     |
 +------------------------------------------+
 |                                          |
 |  +-------------------+  +--------------+ |
 |  |    settings       |  |    users     | |
 |  |  (singleton)      |  | (singleton)  | |
 |  +-------------------+  +--------------+ |
 |  | id (PK, CHECK=1)  |  | id (PK)      | |
 |  | base_url (TEXT)   |  | first_name   | |
 |  | api_key (TEXT)    |  | last_name    | |
 |  | created_at (TEXT) |  | profile_pic  | |
 |  | updated_at (TEXT) |  |  (BLOB)      | |
 |  +-------------------+  +--------------+ |
 |                                          |
 |  +-------------------+  +--------------+ |
 |  | wallet_addresses  |  |schema_version| |
 |  +-------------------+  +--------------+ |
 |  | id (PK AUTO)      |  | version (PK) | |
 |  | address (UNIQUE)  |  +--------------+ |
 |  | label (TEXT)      |                   |
 |  | created_at (TEXT) |                   |
 |  | updated_at (TEXT) |                   |
 |  +-------------------+                   |
 +------------------------------------------+
           ^                    ^
           |                    |
      +----+----+          +----+----+
      |  Iced   |          |  Tauri  |
      | Wallet  |          | Wallet  |
      +---------+          +---------+
     Same password = same encryption key
\end{minted}

% ──────────────────────────────────────────────────────────────
\section{Architecture: a single encrypted database + a single guarded connection}
\label{sec:single-connection-pattern}

We use a single global connection in both wallets:

\begin{minted}[fontsize=\footnotesize]{rust}
static DB_CONN: OnceLock<Mutex<Connection>> = OnceLock::new();
\end{minted}

This is a deliberate trade-off suitable for a GUI application:

\begin{itemize}
  \item it avoids passing a connection through every call,
  \item it ensures access is serialized (SQLite \code{Connection} is not safe to share without synchronization),
  \item and it keeps the API of the module simple --- any function in the module can call \code{get\_connection()} to get a \code{MutexGuard<Connection>}.
\end{itemize}

The \code{OnceLock} guarantees that \code{init\_database} can only succeed once per process. After that, \code{get\_connection()} serves as the narrow waist: every read/write operation acquires the mutex, executes its query, and releases the lock when the guard drops.

% ──────────────────────────────────────────────────────────────
\section{Cargo dependency: \texttt{rusqlite} with \texttt{bundled-sqlcipher}}
\label{sec:cargo-deps}

Both wallets declare the same dependency:

\begin{minted}[fontsize=\footnotesize]{toml}
rusqlite = { version = "0.31", features = ["bundled-sqlcipher"] }
dirs = "5.0"
\end{minted}

The \code{bundled-sqlcipher} feature compiles SQLCipher from source during \code{cargo build}, so no external SQLCipher library installation is needed. The \code{dirs} crate provides cross-platform data directory resolution (e.g., \code{~/Library/Application Support} on macOS, \code{~/.local/share} on Linux, \code{\%APPDATA\%} on Windows).

% ──────────────────────────────────────────────────────────────
\section{Key derivation: deterministic password generation}
\label{sec:key-derivation}

We generate the database password without user interaction. Both wallets use an identical \code{generate\_database\_password} function---in fact, the Tauri version includes the comment ``EXACT replica from bitcoin-wallet-ui-iced'':

\begin{listing}[htbp]
\caption{Deterministic password generation (identical in both wallets)}
\label{lst:ch20-password-gen}
\begin{minted}[fontsize=\footnotesize]{rust}
fn generate_database_password() -> String {
    use std::collections::hash_map::DefaultHasher;
    use std::hash::{Hash, Hasher};

    let mut hasher = DefaultHasher::new();

    // Use username
    if let Ok(username) = std::env::var("USER") {
        username.hash(&mut hasher);
    } else if let Ok(username) = std::env::var("USERNAME") {
        username.hash(&mut hasher);
    }

    // Use home directory
    if let Some(home) = dirs::home_dir() {
        home.to_string_lossy().hash(&mut hasher);
    }

    // Use application name
    "bitcoin-wallet-ui".hash(&mut hasher);

    // Convert to hex string
    format!("{:x}", hasher.finish())
}
\end{minted}
\end{listing}

We emphasize \textbf{zero-interaction UX} with this design: the user is never prompted for a passphrase, but the file is still encrypted at rest. The password is deterministic---the same user on the same machine always produces the same key---which means both the Iced and Tauri wallets can open the same database file.

We use three inputs for the hash:

\begin{enumerate}
  \item \textbf{Username}: from \code{\$USER} (Unix/macOS) or \code{\$USERNAME} (Windows).
  \item \textbf{Home directory}: via \code{dirs::home\_dir()}, which resolves platform-specific paths.
  \item \textbf{Application name}: the literal string \code{"bitcoin-wallet-ui"}.
\end{enumerate}

We use \code{DefaultHasher} to produce a 64-bit hash, which we convert to a hex string. This is not cryptographically strong key derivation (we don't use PBKDF2 or Argon2), but it is sufficient for a local development/learning tool: the encryption prevents casual file inspection, and the determinism avoids lockout.

\begin{notebox}
Because the application name \code{"bitcoin-wallet-ui"} is the same in both Iced and Tauri wallets, and because both store the file at the same path (\code{<data\_dir>/bitcoin-wallet-ui/bitcoin-wallet.db}), the two wallet apps share the same database. This is intentional --- a user who switches between the two frontends sees the same wallets and settings.
\end{notebox}

% ──────────────────────────────────────────────────────────────
\section{Database initialization}
\label{sec:init-database}

The initialization sequence is identical in both wallets. The caller generates the password and passes it in:

\begin{minted}[fontsize=\footnotesize]{rust}
// In main() --- identical in both wallets
let db_password = generate_database_password();
if let Err(e) = database::init_database(&db_password) {
    eprintln!("Failed to initialize database: {}", e);
    // Continue anyway - settings will use defaults
}
\end{minted}

\subsection{Annotated listing: \texttt{init\_database}}
\label{sec:init-database-fn}

\begin{listing}[H]
\caption{\code{init\_database}: establishing the encrypted connection}
\label{lst:ch20-init-db}
\begin{minted}[fontsize=\footnotesize]{rust}
pub fn init_database(password: &str) -> SqliteResult<()> {
    let db_path = get_database_path();

    // 1) Ensure the parent directory exists. This makes the module robust on
    // first run.
    if let Some(parent) = db_path.parent() {
        std::fs::create_dir_all(parent).map_err(|e| {
            // Map an I/O error into a rusqlite error so we keep one error
            // domain.
            rusqlite::Error::SqliteFailure(
                rusqlite::ffi::Error::new(rusqlite::ffi::SQLITE_CANTOPEN),
                Some(format!("Failed to create database directory: {}", e)),
            )
        })?;
    }

    // 2) Open (or create) the SQLite file.
    let conn = Connection::open(&db_path)?;

    // 3) Critical SQLCipher step: set the encryption key *immediately* after
    // opening.
    // If the key is wrong, subsequent reads will fail with "file is
    // encrypted...".
    conn.pragma_update(None, "key", password)?;

    // 4) Enable foreign keys (good SQLite hygiene; harmless if unused).
    conn.pragma_update(None, "foreign_keys", "ON")?;

    // 5) Create base schema + apply migrations before the connection is
    // exposed.
    create_tables(&conn)?;
    run_migrations(&conn)?;

    // 6) Store the connection globally for the rest of the process lifetime.
    DB_CONN.set(Mutex::new(conn)).map_err(|_| {
        rusqlite::Error::SqliteFailure(
            rusqlite::ffi::Error::new(rusqlite::ffi::SQLITE_MISUSE),
            Some("Database already initialized".to_string()),
        )
    })?;

    Ok(())
}
\end{minted}
\end{listing}

We present this six-step sequence because it embodies a pattern we reuse:

\begin{enumerate}
  \item \textbf{Ensure directories exist} --- \code{create\_dir\_all} is idempotent.
  \item \textbf{Open the file} --- \code{Connection::open} creates the file if it doesn't exist.
  \item \textbf{Set the encryption key} --- this must be the very first pragma. If we ran any other SQL first, SQLCipher would treat the database as unencrypted.
  \item \textbf{Enable foreign keys} --- SQLite disables them by default.
  \item \textbf{Create tables and run migrations} --- before exposing the connection to any caller.
  \item \textbf{Store globally} --- \code{OnceLock::set} fails if already initialized, which we surface as a \code{SQLITE\_MISUSE} error.
\end{enumerate}

\subsection{\texttt{get\_database\_path}}
\label{sec:db-path}

The database file location is resolved via the \code{dirs} crate:

\begin{minted}[fontsize=\footnotesize]{rust}
fn get_database_path() -> PathBuf {
    let mut path = dirs::data_dir().unwrap_or_else(|| PathBuf::from("."));
    path.push("bitcoin-wallet-ui");
    path.push(DB_FILENAME);
    path
}
\end{minted}

This resolves to platform-specific paths:

\begin{table}[htbp]
\centering\small
\begin{tabularx}{0.8\textwidth}{lX}
\toprule
\textbf{Platform} & \textbf{Path} \\
\midrule
macOS & \code{\textasciitilde/Library/Application Support/bitcoin-wallet-ui/bitcoin-wallet.db} \\
Linux & \code{\textasciitilde/.local/share/bitcoin-wallet-ui/bitcoin-wallet.db} \\
Windows & \code{\%APPDATA\%/bitcoin-wallet-ui/bitcoin-wallet.db} \\
\bottomrule
\end{tabularx}
\end{table}

% ──────────────────────────────────────────────────────────────
\section{Connection acquisition: \texttt{get\_connection}}
\label{sec:get-connection}

\code{get\_connection} is the narrow waist of the module: every read/write operation goes through it. It ensures the database was initialized and acquires the global mutex:

\begin{listing}[htbp]
\caption{\code{get\_connection}: the singleton access point}
\label{lst:ch20-get-conn}
\begin{minted}[fontsize=\footnotesize]{rust}
fn get_connection() -> SqliteResult<
    std::sync::MutexGuard<'static, Connection>
> {
    DB_CONN
        .get()
        .ok_or_else(|| {
            // The database was never initialized
            // --- UI must call init_database at startup.
            rusqlite::Error::SqliteFailure(
                rusqlite::ffi::Error::new(
                    rusqlite::ffi::SQLITE_MISUSE
                ),
                Some("Database not initialized".to_string()),
            )
        })?
        .lock()
        .map_err(|_| {
            // Poisoned mutex implies a panic
            // while holding the lock.
            rusqlite::Error::SqliteFailure(
                rusqlite::ffi::Error::new(
                    rusqlite::ffi::SQLITE_MISUSE
                ),
                Some("Failed to acquire lock".to_string()),
            )
        })
}
\end{minted}
\end{listing}

Two things can go wrong: \code{OnceLock::get()} returns \code{None} if \code{init\_database} was never called, and \code{Mutex::lock()} returns \code{Err} if a previous caller panicked while holding the lock (poisoning the mutex). Both are mapped into the \code{rusqlite::Error} domain so every database function has a uniform error type.

% ──────────────────────────────────────────────────────────────
\section{Schema and defaults: \texttt{create\_tables}}
\label{sec:create-tables}

We read \code{create\_tables} with three lenses:

\begin{itemize}
  \item \textbf{Singleton settings}: \code{settings.id} has \code{CHECK (id = 1)} so there can be only one row.
  \item \textbf{Uniqueness}: \code{wallet\_addresses.address} is \code{UNIQUE} to prevent duplicates.
  \item \textbf{Schema versioning}: \code{schema\_version} is created and initialized to \code{SCHEMA\_VERSION} (currently 2).
\end{itemize}

\subsection{Settings and wallet addresses tables}
\label{sec:tables-part1}

\begin{listing}[H]
\caption{Settings, wallets, and schema version tables}
\label{lst:ch20-create-tables}
\begin{minted}[fontsize=\footnotesize]{rust}
fn create_tables(conn: &Connection) -> SqliteResult<()> {
    // -- Settings table --
    // CHECK(id = 1) prevents multiple rows (singleton pattern).
    conn.execute(
        "CREATE TABLE IF NOT EXISTS settings (
            id INTEGER PRIMARY KEY CHECK (id = 1),
            base_url TEXT NOT NULL DEFAULT '...',
            api_key TEXT NOT NULL DEFAULT '',
            created_at TEXT NOT NULL DEFAULT (datetime('now')),
            updated_at TEXT NOT NULL DEFAULT (datetime('now'))
        )",
        [],
    )?;

    // Populate singleton row on first run.
    let count: i64 = conn.query_row(
        "SELECT COUNT(*) FROM settings",
        [],
        |row| row.get(0)
    )?;
    if count == 0 {
        conn.execute(
            "INSERT INTO settings (id, base_url, api_key) \
             VALUES (1, 'http://127.0.0.1:8080', ?)",
            params![
                std::env::var("BITCOIN_API_WALLET_KEY")
                    .unwrap_or_else(|_| "wallet-secret".to_string())
            ],
        )?;
    }

    // -- Wallet addresses table --
    conn.execute(
        "CREATE TABLE IF NOT EXISTS wallet_addresses (
            id INTEGER PRIMARY KEY AUTOINCREMENT,
            address TEXT NOT NULL UNIQUE,
            label TEXT,
            created_at TEXT NOT NULL DEFAULT (datetime('now')),
            updated_at TEXT NOT NULL DEFAULT (datetime('now'))
        )",
        [],
    )?;

    conn.execute(
        "CREATE INDEX IF NOT EXISTS \
         idx_wallet_addresses_address \
         ON wallet_addresses(address)",
        [],
    )?;

    // -- User profile table --
    conn.execute(
        "CREATE TABLE IF NOT EXISTS users (
            id INTEGER PRIMARY KEY CHECK (id = 1),
            first_name TEXT,
            last_name TEXT,
            profile_picture BLOB,
            created_at TEXT NOT NULL DEFAULT (datetime('now')),
            updated_at TEXT NOT NULL DEFAULT (datetime('now'))
        )",
        [],
    )?;

    // -- Schema version table --
    conn.execute(
        "CREATE TABLE IF NOT EXISTS schema_version (
            version INTEGER PRIMARY KEY
        )",
        [],
    )?;

    let version_count: i64 = conn.query_row(
        "SELECT COUNT(*) FROM schema_version",
        [],
        |row| row.get(0)
    )?;
    if version_count == 0 {
        conn.execute(
            "INSERT INTO schema_version (version) VALUES (?)",
            params![SCHEMA_VERSION],
        )?;
    }

    Ok(())
}
\end{minted}
\end{listing}

We introduce a pattern here called the \textbf{singleton table}: \code{CHECK (id = 1)} on both \code{settings} and \code{users} prevents more than one row. This is a simple and effective alternative to a key-value store when we have a fixed set of fields.

% ──────────────────────────────────────────────────────────────
\section{Migrations: evolving the database safely}
\label{sec:migrations}

We implement migrations to handle one of the most common real-world problems: shipping a new binary against a database created by an older version.

The current code includes one migration, from schema version 1 to version 2:

\begin{itemize}
  \item \textbf{Change}: \code{users.profile\_picture\_path} (TEXT) \textasciitilde{} \code{users.profile\_picture} (BLOB).
  \item \textbf{Challenge}: SQLite does not support \code{ALTER COLUMN} to change a column type.
  \item \textbf{Solution}: Recreate the table via \code{execute\_batch} within a transaction.
\end{itemize}

\subsection{\texttt{run\_migrations}}
\label{sec:run-migrations}

\begin{listing}[H]
\caption{Migration from schema v1 to v2}
\label{lst:ch20-migrations}
\begin{minted}[fontsize=\footnotesize]{rust}
fn run_migrations(conn: &Connection) -> SqliteResult<()> {
    let current_version: i32 = conn
        .query_row("SELECT version FROM schema_version LIMIT 1", [], |row| {
            row.get(0)
        })
        .unwrap_or(0);

    // Migration 1 -> 2: detect old column and recreate table
    if current_version < 2 {
        let table_info = conn.prepare("PRAGMA table_info(users)")?;
        let has_old_column = table_info.exists()?;

        if has_old_column {
            conn.execute_batch(
                "BEGIN TRANSACTION;
                 CREATE TABLE users_new (
                     id INTEGER PRIMARY KEY CHECK (id = 1),
                     first_name TEXT,
                     last_name TEXT,
                     profile_picture BLOB,
                     created_at TEXT NOT NULL,
                     updated_at TEXT NOT NULL
                 );
                 INSERT INTO users_new
                   SELECT id, first_name, last_name, NULL,
                          created_at, updated_at FROM users;
                 DROP TABLE users;
                 ALTER TABLE users_new RENAME TO users;
                 COMMIT;"
            )?;
        }

        conn.execute("UPDATE schema_version SET version = ?", params![2])?;
    }

    if current_version < SCHEMA_VERSION {
        conn.execute(
            "UPDATE schema_version SET version = ?",
            params![SCHEMA_VERSION]
        )?;
    }

    Ok(())
}
\end{minted}
\end{listing}

We use SQLite's \code{PRAGMA table\_info()} to introspect the current schema before making changes. This makes the migration idempotent: running it against a database that has already been migrated is a no-op. The ``recreate table'' technique---create a new table with the desired schema, copy the data, drop the old table, rename the new table---is the standard approach for SQLite schema changes that go beyond what \code{ALTER TABLE} can express.

% ──────────────────────────────────────────────────────────────
\section{Data types}
\label{sec:db-types}

We define the same two public types in both wallets. The only difference is that the Tauri wallet adds \code{Serialize} and \code{Deserialize} derives (required for Tauri IPC serialization):

\begin{listing}[htbp]
\caption{Public data types (\code{Settings} and \code{WalletAddress})}
\label{lst:ch20-types}
\begin{minted}[fontsize=\footnotesize]{rust}
// Iced wallet
#[derive(Debug, Clone)]
pub struct Settings {
    pub base_url: String,
    pub api_key: String,
}

#[derive(Debug, Clone)]
pub struct WalletAddress {
    pub address: String,
    pub label: Option<String>,
    pub created_at: String,
}

// Tauri wallet (adds Serialize/Deserialize)
#[derive(Debug, Clone, Serialize, Deserialize)]
pub struct Settings {
    pub base_url: String,
    pub api_key: String,
}

impl WalletAddress {
    pub fn new(address: String, label: Option<String>) -> Self {
        Self {
            address,
            label,
            created_at: String::new(), // Set by database
        }
    }
}
\end{minted}
\end{listing}

We intentionally leave \code{created\_at} empty: the database's \code{DEFAULT (datetime('now'))} clause assigns the timestamp on insert.

% ──────────────────────────────────────────────────────────────
\section{Settings persistence: load and save}
\label{sec:settings-crud}

\subsection{\texttt{load\_settings}}
\label{sec:load-settings}

\begin{listing}[H]
\caption{Reading settings from the singleton row}
\label{lst:ch20-load-settings}
\begin{minted}[fontsize=\footnotesize]{rust}
pub fn load_settings() -> SqliteResult<Settings> {
    let conn = get_connection()?;

    conn.query_row(
        "SELECT base_url, api_key FROM settings WHERE id = 1",
        [],
        |row| {
            Ok(Settings {
                base_url: row.get(0)?,
                api_key: row.get(1)?,
            })
        },
    )
}
\end{minted}
\end{listing}

\subsection{\texttt{save\_settings}}
\label{sec:save-settings}

\begin{listing}[H]
\caption{Updating settings}
\label{lst:ch20-save-settings}
\begin{minted}[fontsize=\footnotesize]{rust}
pub fn save_settings(settings: &Settings) -> SqliteResult<()> {
    let conn = get_connection()?;

    conn.execute(
        "UPDATE settings SET base_url = ?, \
         api_key = ?, \
         updated_at = datetime('now') WHERE id = 1",
        params![settings.base_url, settings.api_key],
    )?;

    Ok(())
}
\end{minted}
\end{listing}

% ──────────────────────────────────────────────────────────────
\section{Wallet address persistence: insert-or-update}
\label{sec:wallet-crud}

The wallet list must be stable under these operations:

\begin{itemize}
  \item creating a new wallet (insert),
  \item creating another wallet with the same address (update label; should not duplicate),
  \item loading all wallets (display list, newest first).
\end{itemize}

\subsection{\texttt{save\_wallet\_address} --- upsert pattern}
\label{sec:save-wallet-address}

\begin{listing}[H]
\caption{Insert or update a wallet address}
\label{lst:ch20-save-wallet}
\begin{minted}[fontsize=\footnotesize]{rust}
pub fn save_wallet_address(
    wallet: &WalletAddress
) -> SqliteResult<WalletAddress> {
    let conn = get_connection()?;

    match conn.execute(
        "INSERT INTO wallet_addresses \
         (address, label, updated_at) \
         VALUES (?, ?, datetime('now'))",
        params![wallet.address, wallet.label],
    ) {
        Ok(_) => {
            conn.query_row(
                "SELECT address, label, created_at \
                 FROM wallet_addresses \
                 WHERE address = ?",
                params![wallet.address],
                |row| {
                    Ok(WalletAddress {
                        address: row.get(0)?,
                        label: row.get(1)?,
                        created_at: row.get(2)?,
                    })
                },
            )
        }
        Err(rusqlite::Error::SqliteFailure(err, _))
            if err.code == rusqlite::ErrorCode::ConstraintViolation =>
        {
            conn.execute(
                "UPDATE wallet_addresses SET \
                 label = ?, \
                 updated_at = datetime('now') \
                 WHERE address = ?",
                params![wallet.label, wallet.address],
            )?;

            conn.query_row(
                "SELECT address, label, created_at \
                 FROM wallet_addresses \
                 WHERE address = ?",
                params![wallet.address],
                |row| {
                    Ok(WalletAddress {
                        address: row.get(0)?,
                        label: row.get(1)?,
                        created_at: row.get(2)?,
                    })
                },
            )
        }
        Err(e) => Err(e),
    }
}
\end{minted}
\end{listing}

This provides a stable call pattern: the caller always gets a complete \code{WalletAddress} back, regardless of whether the operation was insert or update. The \code{ConstraintViolation} match arm is the key---it converts what would normally be an error into a fallback update.

\subsection{\texttt{load\_wallet\_addresses}}
\label{sec:load-wallet-addresses}

\begin{listing}[H]
\caption{Loading the entire wallet list}
\label{lst:ch20-load-wallets}
\begin{minted}[fontsize=\footnotesize]{rust}
pub fn load_wallet_addresses() -> SqliteResult<Vec<WalletAddress>> {
    let conn = get_connection()?;
    let mut stmt = conn.prepare(
        "SELECT address, label, created_at \
         FROM wallet_addresses \
         ORDER BY created_at DESC"
    )?;

    let addresses = stmt.query_map([], |row| {
        Ok(WalletAddress {
            address: row.get(0)?,
            label: row.get(1)?,
            created_at: row.get(2)?,
        })
    })?;

    addresses.collect::<Result<Vec<_>, _>>()
}
\end{minted}
\end{listing}

% ──────────────────────────────────────────────────────────────
\section{Additional Tauri operations}
\label{sec:tauri-operations}

The Tauri wallet adds two CRUD operations not in Iced:

\begin{listing}[htbp]
\caption{Tauri-specific deletion and label updates}
\label{lst:ch20-tauri-ops}
\begin{minted}[fontsize=\footnotesize]{rust}
pub fn delete_wallet_address_db(address: &str) -> SqliteResult<()> {
    get_connection()?.execute(
        "DELETE FROM wallet_addresses WHERE address = ?",
        params![address]
    )?;
    Ok(())
}

pub fn update_wallet_label_db(
    address: &str,
    label: &str
) -> SqliteResult<()> {
    get_connection()?.execute(
        "UPDATE wallet_addresses SET \
         label = ?, \
         updated_at = datetime('now') \
         WHERE address = ?",
        params![label, address],
    )?;
    Ok(())
}
\end{minted}
\end{listing}

% ──────────────────────────────────────────────────────────────
\section{Iced vs.\ Tauri: comparison}
\label{sec:iced-vs-tauri}

The database layer is the most shared code between the two wallet frontends. Here is a precise comparison:

\begin{table}[htbp]
\centering\small
\begin{tabularx}{0.95\textwidth}{lXX}
\toprule
\textbf{Aspect} & \textbf{Iced Wallet} & \textbf{Tauri Wallet} \\
\midrule
\textbf{File location} & \code{src/database.rs} & \code{src-tauri/src/database/mod.rs} \\
\textbf{Schema} & 4 tables, v2 & 4 tables, v2 (identical) \\
\textbf{\code{init\_database}} & Identical & Identical \\
\textbf{\code{create\_tables}} & Identical & Identical \\
\textbf{\code{run\_migrations}} & Identical & Identical \\
\textbf{Types} & \code{Debug, Clone} & \code{Debug, Clone, Serialize, Deserialize} \\
\textbf{Save settings} & \code{save\_settings} & \code{save\_settings\_db} \\
\textbf{Delete wallet} & Not implemented & \code{delete\_wallet\_address\_db} \\
\textbf{Update label} & Not implemented & \code{update\_wallet\_label\_db} \\
\textbf{Unit tests} & None & 7 tests in \code{tests.rs} \\
\bottomrule
\end{tabularx}
\end{table}

The takeaway is that the schema and initialization are identical. The Tauri wallet adds \code{Serialize}/\code{Deserialize} (needed for IPC), renames one function to avoid command-layer naming conflicts, and adds delete + label update operations.

% ──────────────────────────────────────────────────────────────
\section{Integration points in the wallet UIs}
\label{sec:integration-points}

This database layer is integrated at three points in each wallet, though the specific mechanism differs.

\subsection{Iced wallet integration}
\label{sec:iced-integration}

\begin{enumerate}
  \item \textbf{Startup} (\code{main.rs}): \code{generate\_database\_password()} \textrightarrow{} \code{database::init\_database(\&db\_password)}.
  \item \textbf{Initial model} (\code{WalletApp::new}): calls \code{load\_settings()} and \code{load\_wallet\_addresses()} to populate the Iced model.
  \item \textbf{Update loop} (\code{update.rs}): \code{Message::SaveSettings} calls \code{save\_settings()}; \code{Message::WalletCreated} calls \code{save\_wallet\_address()}.
\end{enumerate}

Persistence is a side effect in the state machine (\code{update}), not in the view.

\subsection{Tauri wallet integration}
\label{sec:tauri-integration}

\begin{enumerate}
  \item \textbf{Startup} (\code{main.rs}): \code{generate\_database\_password()} \textrightarrow{} \code{database::init\_database(\&db\_password)}.
  \item \textbf{AppState initialization} (\code{config/mod.rs}): \code{AppState::default()} calls \code{load\_settings()} to populate initial state from the database with env var fallback.
  \item \textbf{Tauri commands} (\code{commands/settings.rs}, \code{commands/wallet.rs}): each \code{\#[tauri::command]} handler calls the appropriate database function. For example, \code{create\_wallet} calls the remote API then \code{save\_wallet\_address()} to persist locally.
\end{enumerate}

The Tauri integration is more fine-grained because each command handler is an independent async function, rather than a single \code{update} state machine.

% ──────────────────────────────────────────────────────────────
\section{Unit tests (Tauri wallet)}
\label{sec:unit-tests}

We include a test suite in the Tauri wallet that validates the database layer without SQLCipher encryption (using in-memory SQLite). The test helper \code{setup\_test\_db()} creates the same schema in memory, then each test exercises a specific aspect:

\begin{table}[htbp]
\centering\small
\begin{tabularx}{0.95\textwidth}{lX}
\toprule
\textbf{Test} & \textbf{What it validates} \\
\midrule
\code{test\_schema\_creation} & All four tables exist after setup \\
\code{test\_settings\_crud} & Read defaults, update, verify changes \\
\code{test\_wallet\_address\_crud} & Insert, read, update label, delete \\
\code{test\_wallet\_address\_uniqueness} & Duplicate address insert fails \\
\code{test\_wallet\_list\_ordering} & \code{ORDER BY created\_at DESC} returns newest first \\
\code{test\_schema\_version} & Version table contains expected value (2) \\
\code{test\_settings\_singleton\_constraint} & \code{CHECK(id = 1)} rejects \code{id = 2} \\
\bottomrule
\end{tabularx}
\end{table}

We use \code{Connection::open\_in\_memory()} without encryption, which isolates the database logic from SQLCipher's key management. This is a practical testing pattern: the SQL schema and CRUD logic are the same whether or not the file is encrypted, so in-memory tests cover the logic without the complexity of key setup.

% ──────────────────────────────────────────────────────────────
\section{Summary}
\label{sec:ch20-summary}

\begin{chaptersummary}
\item We integrated SQLCipher for encrypted wallet storage across both Iced and Tauri frontends.
\item We designed database schemas and implemented migrations using the table-recreation technique.
\item We handled thread safety with \code{OnceLock<Mutex<Connection>>} and the singleton table pattern.
\item We implemented deterministic password generation so both wallet UIs share one encrypted database file.
\end{chaptersummary}

In the next chapter, we build the web admin interface --- a React/TypeScript application that provides browser-based access to the blockchain.

\begin{infobox}[Companion Chapter]
Complete source code including both Iced and Tauri implementations, tests, and the shared password generation function is available in the companion code listings appendix. In the print edition, these listings appear in the Appendix: Source Reference.
\end{infobox}

---

\section{Exercises}
\label{sec:ch20-exercises}

\begin{enumerate}
  \item \textbf{Add a New Table} --- Create a new \code{transaction\_memos} table in the SQLCipher database for storing user-provided notes on transactions. Write the migration using the table-recreation technique, implement insert and query functions, and write a test that stores and retrieves a memo.

  \item \textbf{Thread Safety Verification} --- Write a test that spawns two threads, each attempting to write to the database simultaneously through the \code{OnceLock<Mutex<Connection>>} singleton. Verify that both writes succeed without data corruption and that the mutex correctly serializes access.
\end{enumerate}
